\documentclass[11pt,a4paper]{article}

\usepackage{lmodern}
\usepackage[T1]{fontenc}
\usepackage[utf8]{inputenc}
\usepackage[margin=2.6cm]{geometry}
\usepackage{amsmath}
\usepackage{amssymb}
\usepackage{booktabs}
\usepackage{microtype}
\usepackage[round]{natbib}
\usepackage[hidelinks]{hyperref}

\newcommand{\code}[1]{\texttt{#1}}

% A numbered, ruled algorithm block in base LaTeX. The `algorithm` and
% `algorithmicx` packages are not in the TeX Live subset the paper workflow
% installs (texlive-latex-extra ships neither; they sit in texlive-science),
% so the environment is defined here rather than pulled in.
\newcounter{algo}
\newenvironment{algobox}[2]{%
  \par\addvspace{\medskipamount}%
  \refstepcounter{algo}\label{#2}%
  \noindent\rule{\linewidth}{0.8pt}\par\nobreak\vspace{2pt}%
  \noindent\textbf{Algorithm \thealgo.} #1\par\nobreak\vspace{2pt}%
  \noindent\rule{\linewidth}{0.4pt}\par\nobreak\vspace{2pt}%
  \begin{list}{\arabic{algoline}:}{%
    \usecounter{algoline}%
    \setlength{\leftmargin}{2.2em}\setlength{\labelwidth}{1.6em}%
    \setlength{\itemsep}{1pt}\setlength{\parsep}{0pt}\setlength{\topsep}{2pt}}%
}{%
  \end{list}\nobreak\vspace{2pt}%
  \noindent\rule{\linewidth}{0.8pt}\par\addvspace{\medskipamount}%
}
\newcounter{algoline}
\DeclareMathOperator*{\argmin}{arg\,min}
\newcommand{\simplex}{\Delta}
\newcommand{\R}{\mathbb{R}}

\title{\code{proximal-lq}: Projected Gradient Descent for
Simplex-Constrained Linear Least Squares}

\author{Thomas Schmelzer\thanks{\url{https://github.com/tschm/proximal}}}

\date{\today}

\begin{document}

\maketitle

\begin{abstract}
\code{proximal-lq} solves $\min_{x \in \simplex} \tfrac{1}{2}\|Ax-b\|_2^2$,
where $\simplex$ is the probability simplex, with projected gradient descent and
the exact $O(n \log n)$ simplex projection of \citet{duchi2008}. The package is
two functions and roughly two hundred lines of NumPy; this paper states what
they compute, the step size and stopping rule they actually use, and how they
behave on the problems in the test suite. Three things are worth recording. The
implementation takes $\tau = 1/(2L)$ rather than the $\tau = 1/L$ its own
development notes specify --- both are convergent, but the shipped constant is
half the classical one and costs roughly a factor of two in iterations. The
parameter named \code{eps\_rel} is compared against an absolute quantity, the
norm of the last step, so tolerances do not scale with problem size. And the
portfolio fixture the test suite validates against solves a least-squares fit of
$\Sigma x \approx \mathbf{1}$, which we show is a different portfolio from the
long-only minimum-variance one --- variance $0.3087$ against $0.0421$ on the
same covariance matrix --- a distinction easy to lose when the two are described
with the same words. The solver itself converges in tens of iterations and
reproduces the reference solution to $4.4 \times 10^{-6}$ in $\ell_1$.
\end{abstract}

\section{Introduction}

Many estimation problems are least squares with the answer constrained to be a
set of proportions: non-negative, summing to one. Long-only fully-invested
portfolio weights \citep{markowitz1952} are the canonical example; mixture
weights, non-negative sparse codes and discrete probability estimates are the
same shape. Formally, with $A \in \R^{m \times n}$ and $b \in \R^m$,
\begin{equation}
  \min_{x \in \R^n} \; f(x) = \tfrac{1}{2} \| A x - b \|_2^2
  \quad \text{subject to} \quad x \in \simplex,
  \qquad
  \simplex = \Bigl\{ x \in \R^n : x \ge 0, \; \textstyle\sum_i x_i = 1 \Bigr\}.
  \label{eq:problem}
\end{equation}

This is a convex quadratic program, so a general-purpose solver
\citep{diamond2016} will handle it. The reason not to use one is that
Eq.~\eqref{eq:problem} has structure a general solver cannot exploit: the
constraint set is simple enough that projecting onto it is cheaper than one
gradient evaluation. Projection onto $\simplex$ costs $O(n \log n)$
\citep{duchi2008}, against $O(n^2)$ for a single gradient once the Gram matrix
is formed. When the projection is the cheap step, projected gradient descent is
the natural method: no factorisations, no interior point, no KKT system, and an
iterate that is feasible after the first step and stays feasible.

\code{proximal-lq} is that method, kept deliberately small --- two public
functions, NumPy \citep{harris2020} as the only dependency. This paper documents
it precisely enough to be used with confidence, which includes being explicit
about the three places where the shipped behaviour differs from what a reader of
the documentation would assume.

\section{Algorithm}

\subsection{Proximal gradient descent}

Write Eq.~\eqref{eq:problem} as the sum of a smooth part and an indicator,
\begin{equation}
  \min_x \; f(x) + g(x),
  \qquad
  f(x) = \tfrac{1}{2}\|Ax-b\|_2^2,
  \qquad
  g(x) = \iota_{\simplex}(x) =
  \begin{cases}
    0 & x \in \simplex \\
    +\infty & \text{otherwise.}
  \end{cases}
\end{equation}
The proximal gradient method --- forward--backward splitting
\citep{parikh2014} --- iterates
\begin{equation}
  x^{k+1} = \operatorname{prox}_{\tau g}\bigl( x^k - \tau \nabla f(x^k) \bigr).
  \label{eq:pg}
\end{equation}
For an indicator function the proximal operator is the Euclidean projection, so
Eq.~\eqref{eq:pg} is exactly projected gradient descent,
$x^{k+1} = P_{\simplex}(x^k - \tau \nabla f(x^k))$.

The gradient is
\begin{equation}
  \nabla f(x) = A^{\top}(Ax - b) = G x - c,
  \qquad G = A^{\top} A, \quad c = A^{\top} b,
\end{equation}
and $G$ and $c$ are formed once, outside the loop. Each iteration then costs one
matrix--vector product, $O(n^2)$, plus one projection, $O(n \log n)$ --- the
problem's $m$ enters only in the $O(mn^2)$ setup.

\subsection{Step size}

$\nabla f$ is Lipschitz with constant $L = \|A^{\top}A\|_2 = \sigma_{\max}(A)^2$,
the spectral norm of the Gram matrix. Projected gradient descent converges for
any fixed $\tau \in (0, 2/L)$, and the textbook choice $\tau = 1/L$ gives
\begin{equation}
  f(x^k) - f(x^{\star}) \;\le\; \frac{\|x^0 - x^{\star}\|_2^2}{2 \tau k}
  \label{eq:rate}
\end{equation}
--- the $O(1/k)$ rate \citep{nesterov2004,beck2009}.

The implementation uses
\begin{equation}
  \tau = \frac{1}{2L},
  \label{eq:step}
\end{equation}
half the classical value, falling back to $\tau = 1$ when $L < 10^{-15}$ so that
an all-but-zero $A$ does not divide by zero. This is a discrepancy worth
recording: the package's development notes state $\tau = 1/L$, while the code
takes $0.5/L$. Both lie inside the convergent range, so nothing is broken, but
by Eq.~\eqref{eq:rate} the shipped constant roughly doubles the iteration count
for a given accuracy --- a conservative choice, not a wrong one, and the numbers
in Section~\ref{sec:results} are the numbers it produces.

\subsection{Projection onto the simplex}

The inner problem is
\begin{equation}
  P_{\simplex}(v) = \argmin_{x \ge 0, \, \sum_i x_i = s} \tfrac{1}{2}\|x - v\|_2^2 .
\end{equation}
Its solution is a soft threshold: $x_i = \max(v_i - \theta, 0)$ for the unique
$\theta$ making the sum $s$. The KKT conditions give this immediately --- the
multiplier on the equality constraint is the same scalar $\theta$ for every
coordinate, and the non-negativity multipliers are active exactly where
$v_i \le \theta$. The only work is finding $\theta$, and sorting makes it
explicit.

\begin{algobox}{Simplex projection \citep{duchi2008}}{alg:proj}
\item[] \textbf{Input:} $v \in \R^n$, radius $s > 0$.
\item $\mu \gets \operatorname{sort}(v)$, in decreasing order.
\item $\theta_j \gets \bigl(\sum_{i \le j} \mu_i - s\bigr) / j$ for each $j$ ---
      the running mean of the top $j$ entries, offset by $s$.
\item $\rho \gets \max \{ j : \mu_j > \theta_j \}$.
\item Return $x_i = \max(v_i - \theta_\rho, 0)$ for each $i$.
\end{algobox}

Algorithm~\ref{alg:proj} is exact --- not an approximation or an inner iteration
--- and costs $O(n \log n)$, dominated by the sort, with $O(n)$ extra space. The
implementation is a four-line vectorised transcription. Linear-time variants
exist \citep{condat2016}; at the sizes this package targets the sort is not the
bottleneck.

Two properties inherited from exactness are worth naming, because the package's
property tests assert them: the projection is idempotent, and it is the identity
on points already in $\simplex$. A radius $s \neq 1$ is supported via the
\code{rad} argument, which makes the same routine a projection onto a scaled
simplex.

\subsection{Stopping rule}

The iteration halts when the step becomes small,
\begin{equation}
  \| x^{k} - x^{k+1} \|_2 \;\le\; \varepsilon,
  \label{eq:stop}
\end{equation}
or when a cap of \code{max\_iter} iterations is reached. For a projected
gradient iteration this is a defensible surrogate for optimality: $x$ is optimal
precisely when it is a fixed point of Eq.~\eqref{eq:pg}, so a vanishing step
means a vanishing fixed-point residual.

Two caveats. First, the argument is named \code{eps\_rel} and documented as a
``relative stopping tolerance'', but Eq.~\eqref{eq:stop} compares an absolute
norm: nothing divides by $\|x^k\|$, by $n$, or by the objective. A tolerance
therefore does not transfer between problems of different scale, and on a
large-$n$ problem the same $\varepsilon$ is a stricter test per coordinate.
Second, hitting \code{max\_iter} is silent --- the last iterate is returned with
no flag, warning, or residual, so a caller who needs to know whether
Eq.~\eqref{eq:stop} was met must re-measure it.

\subsection{Initialisation}

The primal variable starts from a standard normal draw, seeded via the
\code{seed} argument. The draw is not in $\simplex$ --- it is not even
non-negative --- which is harmless, because the first projection in
Eq.~\eqref{eq:pg} maps it into $\simplex$ and every subsequent iterate stays
there. The problem is convex, so the optimum does not depend on the seed; the
seed exists to make the iteration count and any tie-breaking reproducible.

\section{Numerical results}
\label{sec:results}

\subsection{Validation against a reference solution}

The test suite pins the solver against the covariance matrix from the
critical-line algorithm dataset of \citet{bailey2013} --- ten assets, condition
number $9.21$ --- solving Eq.~\eqref{eq:problem} with $A = \Sigma$ and
$b = \mathbf{1}$. The reference weight vector is reproduced to
$\|x - x_{\mathrm{ref}}\|_1 = 4.4 \times 10^{-6}$ at the default tolerance, with
seven of the ten weights driven to exactly zero.

\begin{table}[t]
\centering
\caption{Iterations to satisfy Eq.~\eqref{eq:stop} on the ten-asset reference
problem, at $\tau = 1/(2L)$.}
\label{tab:iters}
\begin{tabular}{lrr}
\toprule
$\varepsilon$ & Iterations & Final step norm \\
\midrule
$10^{-4}$ & $23$ & $8.69 \times 10^{-5}$ \\
$10^{-6}$ & $36$ & $8.14 \times 10^{-7}$ \\
$10^{-8}$ & $49$ & $7.62 \times 10^{-9}$ \\
\bottomrule
\end{tabular}
\end{table}

Table~\ref{tab:iters} shows the iteration count growing roughly linearly in the
number of digits requested --- about $13$ iterations per two decades --- which is
the geometric behaviour expected of a fixed-step projected gradient on a
strongly convex problem, better than the $O(1/k)$ worst case of
Eq.~\eqref{eq:rate}. The absolute counts are small: the whole solve is tens of
matrix--vector products.

\subsection{Scaling and sparsity}

\begin{table}[t]
\centering
\caption{Random dense problems, $A \in \R^{2n \times n}$ with i.i.d.\ standard
normal entries, default settings. \emph{Non-zeros} counts weights above
$10^{-9}$.}
\label{tab:scale}
\begin{tabular}{lrrr}
\toprule
$n$ & Wall time & $\sum_i x_i$ & Non-zeros \\
\midrule
$10$   & $1.1$ ms  & $1.0000000000$ & $7$  \\
$100$  & $1.8$ ms  & $1.0000000000$ & $29$ \\
$1000$ & $51.5$ ms & $1.0000000000$ & $91$ \\
\bottomrule
\end{tabular}
\end{table}

Table~\ref{tab:scale} reports single-threaded timings on a laptop; they are
indicative, not a benchmark. Two things in it matter more than the times.

The sum constraint holds to ten decimal places at every size, which is a
consequence of the projection being exact rather than of the tolerance: every
iterate is feasible, so feasibility does not degrade with the number of
iterations the way it does for a penalty or an augmented-Lagrangian scheme.

The solutions are sparse, and increasingly so in relative terms: $91$ non-zeros
out of $1000$ variables. This is intrinsic to the geometry, not a regulariser
someone added. The unconstrained least-squares solution of a random
over-determined system has no zeros; projecting onto the simplex thresholds
every coordinate below $\theta$ to exactly zero, and the optimum of a linear
objective over a polytope sits at a face. Users should expect exact zeros, and
should not read them as evidence of a sparsity penalty.

\subsection{What the portfolio fixture is not}
\label{sec:notgmv}

The reference problem of Section~\ref{sec:results} substitutes a covariance
matrix for $A$ and $\mathbf{1}$ for $b$, so it minimises
$\tfrac{1}{2}\|\Sigma x - \mathbf{1}\|_2^2$ over the simplex. It is tempting ---
and both the package's example script and its README invite the reading --- to
call the result an optimal portfolio without qualification. It is worth being
precise about which portfolio, because it is not the one most readers will
assume.

\begin{table}[t]
\centering
\caption{Two long-only fully-invested portfolios on the same ten-asset
covariance matrix. Both are computed with this package's projection.}
\label{tab:gmv}
\begin{tabular}{lrr}
\toprule
& Objective solved & Variance $x^{\top}\Sigma x$ \\
\midrule
Least-squares fixture & $\tfrac{1}{2}\|\Sigma x - \mathbf{1}\|_2^2$ & $0.30865$ \\
Minimum variance      & $\tfrac{1}{2} x^{\top}\Sigma x$             & $0.04212$ \\
\bottomrule
\end{tabular}
\end{table}

Table~\ref{tab:gmv} puts the two side by side. The least-squares fixture yields a
portfolio whose variance is more than seven times that of the minimum-variance
portfolio on the same inputs, and the two weight vectors are far apart
($\ell_1$ distance $1.81$ out of a maximum of $2$). Solving $\Sigma x \approx
\mathbf{1}$ in least squares recovers a minimum-variance-like direction only
when the unconstrained solution $\Sigma^{-1}\mathbf{1}$ happens to be feasible
for the simplex --- that is, non-negative --- and here it is not, at which point
the least-squares residual and the variance are simply different objectives with
different minimisers.

Nothing about the solver is at fault; this is a statement about the fixture. Note
also that minimum variance is itself in reach of this package: with $A$ any
matrix satisfying $A^{\top}A = \Sigma$ --- a Cholesky or symmetric square-root
factor --- and $b = 0$, Eq.~\eqref{eq:problem} \emph{is} the long-only
minimum-variance problem. The second row of Table~\ref{tab:gmv} was produced by
iterating this package's own projection on $\nabla f(x) = \Sigma x$.

\section{Limitations}

\paragraph{No acceleration.} The iteration is plain proximal gradient. Nesterov
momentum \citep{beck2009} would improve the worst-case rate from $O(1/k)$ to
$O(1/k^2)$ at the cost of one extra vector of state, and the conservative step
of Eq.~\eqref{eq:step} leaves further headroom: a backtracking line search would
adapt $\tau$ to the local curvature instead of to a global bound. Neither is
implemented.

\paragraph{No convergence certificate.} The solver returns an array. It does not
return the achieved residual, the iteration count, or whether
Eq.~\eqref{eq:stop} was ever satisfied, and \code{max\_iter} is reached
silently. For an automated pipeline this is the most consequential gap in the
API: a non-converged answer is indistinguishable from a converged one.

\paragraph{Dense and in-memory.} $G = A^{\top}A$ is formed explicitly, which is
$O(n^2)$ storage and $O(mn^2)$ work regardless of any sparsity in $A$. For large
sparse $A$ it would be cheaper to apply $A$ and $A^{\top}$ separately and never
form $G$.

\paragraph{Fixed problem shape.} Only Eq.~\eqref{eq:problem} is supported: one
matrix, one vector, one simplex. There is no weighted norm, no regularisation
term, no box constraints beyond the simplex, no equality constraints other than
the sum, and no warm start --- each call begins from a fresh random point, so a
sequence of related problems cannot reuse work.

\section{Conclusion}

Eq.~\eqref{eq:problem} is a problem where the specialised method is both simpler
and faster than the general one, because the projection is exact and cheap. Two
hundred lines of NumPy solve the ten-asset reference problem to $10^{-8}$ in $49$
iterations and a thousand-variable random problem in $51$ ms, with feasibility
holding to ten decimals at every iterate.

Three things a user should carry away, all of them documentation issues rather
than defects. The step size is $1/(2L)$, half the value the development notes
state, so iteration counts are about twice the textbook figure. The tolerance
called \code{eps\_rel} is absolute, so it does not transfer across problem
scales. And the portfolio example solves a least-squares fit, not a
minimum-variance problem --- the same solver will do minimum variance, but only
if it is handed a square-root factor of the covariance matrix and $b=0$.

\appendix

\section{Reproducibility}
\label{app:repro}

All numbers were produced with \code{proximal-lq} 0.2.2 and the covariance
matrix at \path{tests/proximal_lq/resources/CLA_Data.csv} (its first three
rows are means and bounds and are skipped).

\begin{verbatim}
import numpy as np
from proximal_lq import prox_gradient, proj_simplex

S = np.genfromtxt("tests/proximal_lq/resources/CLA_Data.csv",
                  delimiter=",", skip_header=1)[3:]

x = prox_gradient(S, np.ones(10), seed=0)          # the fixture
print(x.round(5), x @ S @ x)

# minimum variance over the same simplex, via this package's projection
L = np.linalg.norm(S, 2)
y = np.ones(10) / 10
for _ in range(200_000):
    z = proj_simplex(y - (0.5 / L) * (S @ y))
    if np.linalg.norm(z - y) < 1e-14:
        y = z
        break
    y = z
print(y.round(5), y @ S @ y)
\end{verbatim}

Iteration counts in Table~\ref{tab:iters} come from an instrumented copy of the
library's loop --- same Gram matrix, same step $0.5/L$, same projection ---
counting passes until Eq.~\eqref{eq:stop} holds. Timings in
Table~\ref{tab:scale} are single runs of \code{prox\_gradient} on
\code{rng.standard\_normal((2*n, n))} with \code{seed=0}, measured with
\code{time.perf\_counter}.

\bibliographystyle{plainnat}
\bibliography{references}

\end{document}
